%=============================================================================
% DEEP ANALYSIS: PSI-FIELD COMPUTING --- SUB-LANDAUER THERMODYNAMICS,
% ALU ARCHITECTURE, RADIATION HARDNESS, AND NEURAL PROCESSING
%
% Patent #10 Deep Dive: Field Computing and Logic Devices
% Author: Gary Thomas Alcock
% Density Field Dynamics Research Programme
%=============================================================================

\documentclass[11pt,twocolumn]{article}

%--- Packages ----------------------------------------------------------------
\usepackage[utf8]{inputenc}
\usepackage[T1]{fontenc}
\usepackage{amsmath,amssymb,amsfonts,amsthm}
\usepackage{mathtools}
\usepackage{physics}
\usepackage{siunitx}
\usepackage{graphicx}
\usepackage{float}
\usepackage{booktabs}
\usepackage{array}
\usepackage{multirow}
\usepackage{enumitem}
\usepackage[colorlinks=true,linkcolor=blue!60!black,citecolor=green!50!black,
            urlcolor=blue!70!black]{hyperref}
\usepackage[margin=1.8cm]{geometry}
\usepackage{fancyhdr}
\usepackage{tikz}
\usetikzlibrary{arrows.meta,calc,patterns,decorations.pathmorphing,
                positioning,shapes.geometric,circuits.logic.US}
\usepackage{pgfplots}
\pgfplotsset{compat=1.18}
\usepackage{caption}
\usepackage{subcaption}
\usepackage{xcolor}
\usepackage{bm}
\usepackage{appendix}

%--- Theorem environments ----------------------------------------------------
\newtheorem{theorem}{Theorem}[section]
\newtheorem{proposition}[theorem]{Proposition}
\newtheorem{lemma}[theorem]{Lemma}
\newtheorem{corollary}[theorem]{Corollary}
\theoremstyle{definition}
\newtheorem{definition}[theorem]{Definition}
\theoremstyle{remark}
\newtheorem{remark}[theorem]{Remark}

%--- Custom commands ---------------------------------------------------------
\newcommand{\psifield}{\psi}
\newcommand{\DFD}{DFD}
\newcommand{\kB}{k_{\mathrm{B}}}
\newcommand{\EL}{E_{\mathrm{L}}}
\newcommand{\Qpsi}{Q_\psi}
\newcommand{\astar}{a_*}
\newcommand{\Wfunc}{W}
\DeclareMathOperator{\sgn}{sgn}
\DeclareMathOperator{\sech}{sech}
\DeclareMathOperator{\Tr}{Tr}

%--- Headers -----------------------------------------------------------------
\pagestyle{fancy}
\fancyhf{}
\fancyhead[L]{\small\textit{Deep Analysis: $\psi$-Field Computing}}
\fancyhead[R]{\small\textit{G.\,T.\,Alcock}}
\fancyfoot[C]{\thepage}
\renewcommand{\headrulewidth}{0.4pt}

%=============================================================================
\begin{document}

\twocolumn[
\begin{@twocolumnfalse}
\begin{center}
{\LARGE\bfseries Deep Analysis of $\psi$-Field Computing:\\[4pt]
Sub-Landauer Thermodynamics, ALU Architecture,\\[4pt]
Radiation Hardness, and Neural Processing}

\vspace{1em}

{\large Gary Thomas Alcock}\\[4pt]
{\small Density Field Dynamics Research Programme}\\[2pt]
{\small\texttt{densityfielddynamics.org}}

\vspace{1em}
{\small\today}

\vspace{1.5em}

\begin{minipage}{0.92\textwidth}
\small\textbf{Abstract.}
We present a rigorous deep analysis of computation using the scalar
refractive field $\psi$ of Density Field Dynamics. Five major results are
established. \textbf{(1)}~Starting from the Jarzynski equality and
Crooks fluctuation theorem, we prove that reversible $\psi$-field logic
with cavity quality factor $Q > \pi\alpha_{\mathrm{rel}}/\ln 2 \approx
182$ dissipates strictly less energy per operation than the Landauer
bound $\kB T \ln 2$, with the exact entropy production per operation
given by $\Delta S_{\mathrm{op}} = \pi \alpha_{\mathrm{rel}} \kB /
Q$. We show this does not violate the second law because reversible
gates perform no information erasure. \textbf{(2)}~We compare
$\psi$-field computing quantitatively against 12 experimental and
theoretical benchmarks spanning Landauer-limit tests, reversible
computing proposals, adiabatic CMOS, superconducting logic (SFQ, AQFP),
photonic computing, and neuromorphic platforms. \textbf{(3)}~We design a
complete 4-bit arithmetic logic unit (ALU) from $\psi$-field gates,
computing energy, latency, and throughput for all standard operations and
comparing against 7\,nm CMOS, AQFP, and photonic equivalents.
\textbf{(4)}~We derive the single-event upset cross-section for
$\psi$-field memory from first principles, demonstrating exponential
suppression with the number of participating cavity modes, and compute
numerical values in $\mathrm{cm}^2/\mathrm{bit}$ for comparison against
SRAM and MRAM. \textbf{(5)}~We design a complete $\psi$-field neural
network layer with explicit forward and backward pass equations, show how
gradient descent maps onto $\psi$-field dynamics, and compute
theoretical TOPS/W figures for comparison against NVIDIA H100, Google TPU
v5, and Intel Loihi~2.
\end{minipage}
\end{center}
\vspace{1.5em}
\end{@twocolumnfalse}
]

%=============================================================================
\section{Introduction}\label{sec:intro}
%=============================================================================

The $\psi$-field computing paradigm, introduced in the companion
paper~\cite{Alcock2025psiComputing}, establishes that the scalar
refractive field of Density Field Dynamics provides a natural substrate
for energy-efficient, reversible computation. That work developed the
basic gate set, memory architecture, and engineering demonstrator design.

This paper goes substantially deeper on five fronts: (i)~a
thermodynamically rigorous proof that $\psi$-field reversible logic
beats the Landauer bound, grounded in the Jarzynski equality and
fluctuation theorems; (ii)~a comprehensive quantitative comparison
against 12 competing technologies; (iii)~a complete 4-bit ALU design
with full performance analysis; (iv)~a first-principles derivation of
radiation hardness; and (v)~a complete neural network layer design with
training dynamics.

Throughout, we use the DFD formalism of~\cite{Alcock2025DFD}: the
single scalar field $\psi(\mathbf{x},t)$ defines a refractive index
$n = e^\psi$ on flat Minkowski spacetime, with field equation
\begin{equation}
\nabla\cdot\!\left[\mu\!\left(\frac{|\nabla\psi|}{\astar}\right)
\nabla\psi\right] = \frac{4\pi G}{c^2}(\rho - \bar\rho),
\label{eq:field_eq}
\end{equation}
where $\mu(x) = x/(1+x)$ is the simple interpolating function derived
from $S^3$ microsector topology.

%=============================================================================
\section{Rigorous Sub-Landauer Energy Theorem}\label{sec:landauer}
%=============================================================================

\subsection{Landauer's Principle: Statement and Scope}\label{sec:landauer_statement}

Landauer's principle~\cite{Landauer1961} states that any logically
irreversible computation---one that maps multiple distinct input states
to a single output state---must dissipate at least
\begin{equation}
\EL = \kB T \ln 2
\label{eq:landauer_bound}
\end{equation}
per bit of information erased, where $T$ is the temperature of the
thermal reservoir. At $T = 300\,\mathrm{K}$, $\EL \approx 2.87 \times
10^{-21}\,\mathrm{J}$; at $T = 1.5\,\mathrm{K}$, $\EL \approx 1.43
\times 10^{-23}\,\mathrm{J}$.

The key qualifier is \emph{logically irreversible}. A logically
reversible operation---a bijective map on the state space---erases no
information and is therefore not bound by Landauer's limit. This
distinction, emphasized by Bennett~\cite{Bennett1973,Bennett1982}, is
the foundation of reversible computing.

\subsection{Fluctuation Theorems: The Modern Framework}\label{sec:fluctuation}

The modern understanding of Landauer's principle rests on the
fluctuation theorems of nonequilibrium statistical mechanics. We build
the proof on two pillars.

\subsubsection{The Jarzynski Equality}

For a system driven from equilibrium state $A$ to state $B$ by an
external protocol $\lambda(t)$ over time interval $[0,\tau]$, the
Jarzynski equality~\cite{Jarzynski1997} states:
\begin{equation}
\langle e^{-W/\kB T}\rangle = e^{-\Delta F/\kB T},
\label{eq:jarzynski}
\end{equation}
where $W$ is the work performed on the system along a particular
trajectory, $\Delta F = F_B - F_A$ is the free energy difference between
the equilibrium states, and $\langle\cdot\rangle$ denotes an average
over all possible trajectories consistent with the protocol.

By Jensen's inequality ($\langle e^{-x}\rangle \geq e^{-\langle
x\rangle}$), this immediately gives the second-law inequality:
\begin{equation}
\langle W\rangle \geq \Delta F.
\label{eq:second_law}
\end{equation}

The dissipated work is
\begin{equation}
W_{\mathrm{diss}} \equiv \langle W\rangle - \Delta F \geq 0.
\label{eq:wdiss}
\end{equation}

\subsubsection{The Crooks Fluctuation Theorem}

Crooks~\cite{Crooks1999} established a stronger result relating the
work distributions of forward and time-reversed protocols:
\begin{equation}
\frac{P_F(W)}{P_R(-W)} = e^{(W - \Delta F)/\kB T},
\label{eq:crooks}
\end{equation}
where $P_F(W)$ is the probability of observing work $W$ in the forward
process and $P_R(-W)$ is the probability of observing work $-W$ in the
reverse process.

\subsection{Application to $\psi$-Field Logic Operations}\label{sec:psi_jarzynski}

We now apply the fluctuation theorem framework to $\psi$-field logic
operations in a cavity with quality factor $Q$.

\subsubsection{System Definition}

Consider a single $\psi$-field cavity mode with Hamiltonian
\begin{equation}
H(q,p;\lambda) = \frac{p^2}{2M_\psi} + V_{\mathrm{eff}}(q;\lambda),
\label{eq:mode_ham}
\end{equation}
where $q$ is the mode amplitude, $p = M_\psi\dot{q}$ is the conjugate
momentum, and $\lambda(t)$ represents the time-dependent control
parameter (bias field, coupling strength, or boundary condition). The
system is weakly coupled to a thermal bath at temperature $T$ through
cavity wall losses, with damping rate
\begin{equation}
\gamma = \frac{\omega_0}{2Q}.
\label{eq:damping}
\end{equation}

The Langevin equation of motion is
\begin{equation}
M_\psi\ddot{q} + 2M_\psi\gamma\dot{q} + \frac{\partial V_{\mathrm{eff}}}{\partial q} = \xi(t),
\label{eq:langevin}
\end{equation}
where $\xi(t)$ is Gaussian white noise with
$\langle\xi(t)\xi(t')\rangle = 4M_\psi\gamma\kB T\,\delta(t-t')$,
satisfying the fluctuation-dissipation theorem.

\subsubsection{Free Energy Analysis for Reversible Gates}

For a logically reversible gate (NOT, SWAP, Toffoli, Fredkin), the
operation is a bijective map on the logical state space. The protocol
$\lambda(t)$ drives the system from state $A$ at $t = 0$ to state $B$
at $t = \tau$, then returns the control parameter to its initial value.

\begin{lemma}[Zero Free Energy Change for Reversible Gates]
\label{lem:zero_dF}
For a reversible $\psi$-field logic gate that returns the control
parameter to its initial value after the operation, the equilibrium free
energy change is
\begin{equation}
\Delta F = 0.
\end{equation}
\end{lemma}

\begin{proof}
The free energy is a state function determined by the Hamiltonian
parameters. Since the control parameter $\lambda$ returns to its initial
value, $H(q,p;\lambda(\tau)) = H(q,p;\lambda(0))$, and therefore
$F(\tau) = F(0)$. The system has been driven from one potential minimum
to another of equal depth (the symmetric double well has degenerate
minima), so $\Delta F = 0$.
\end{proof}

\subsubsection{Dissipated Work Calculation}

\begin{theorem}[Exact Entropy Production per $\psi$-Field Operation]
\label{thm:entropy_production}
The mean dissipated work per reversible logic operation in a $\psi$-field
cavity with quality factor $Q$, operating at temperature $T$ with
reliability parameter $\alpha_{\mathrm{rel}}$ (barrier height $\Delta V
= \alpha_{\mathrm{rel}}\kB T$), is
\begin{equation}
\boxed{
\langle W_{\mathrm{diss}}\rangle
= \frac{\pi\alpha_{\mathrm{rel}}\,\kB T}{Q}
}
\label{eq:exact_dissipation}
\end{equation}
and the corresponding entropy production is
\begin{equation}
\boxed{
\Delta S_{\mathrm{op}} = \frac{\pi\alpha_{\mathrm{rel}}\,\kB}{Q}
}
\label{eq:exact_entropy}
\end{equation}
per operation.
\end{theorem}

\begin{proof}
We proceed in four steps.

\textbf{Step 1: Work decomposition.}
The work performed on the system during one switching operation
(half-period of oscillation, duration $\tau = \pi/\omega_0$) is
\begin{equation}
W = \int_0^\tau \frac{\partial H}{\partial\lambda}\dot{\lambda}\,dt
= \int_0^\tau \frac{\partial V_{\mathrm{eff}}}{\partial\lambda}\dot{\lambda}\,dt.
\end{equation}

\textbf{Step 2: Energy balance.}
During the switching transient, the mode traverses from one potential
minimum ($q = -q_0$) over the barrier ($q = 0$) to the other minimum
($q = +q_0$). The kinetic energy at the barrier top equals the barrier
height: $\frac{1}{2}M_\psi\dot{q}_{\max}^2 = \Delta V$. The total
mechanical energy exchanged during the transient is $E_{\mathrm{mech}} =
\Delta V$ (borrowed during ascent, returned during descent to the
symmetric minimum).

\textbf{Step 3: Friction integral.}
The energy dissipated to the thermal bath is the work done against the
friction force:
\begin{equation}
W_{\mathrm{diss}} = \int_0^\tau 2M_\psi\gamma\,\dot{q}^2\,dt.
\label{eq:friction_integral}
\end{equation}

For the switching trajectory over a symmetric quartic barrier $V(q) =
\beta(q^2 - q_0^2)^2/4$, the equation of motion in the underdamped
limit ($\gamma \ll \omega_0$) gives a trajectory well-approximated by
\begin{equation}
q(t) \approx q_0\tanh\!\left(\frac{\omega_0(t - \tau/2)}{\sqrt{2}}\right),
\end{equation}
with velocity
\begin{equation}
\dot{q}(t) \approx \frac{q_0\omega_0}{\sqrt{2}}\,\sech^2\!\left(\frac{\omega_0(t-\tau/2)}{\sqrt{2}}\right).
\end{equation}

Computing the friction integral~\eqref{eq:friction_integral}:
\begin{align}
W_{\mathrm{diss}} &= 2M_\psi\gamma \cdot \frac{q_0^2\omega_0^2}{2}
\int_0^{\pi/\omega_0}\sech^4\!\left(\frac{\omega_0 t'}{\sqrt{2}}\right)dt'
\nonumber\\
&= M_\psi\gamma q_0^2\omega_0^2
\cdot\frac{\sqrt{2}}{\omega_0}\cdot\frac{4}{3}\nonumber\\
&= \frac{4\sqrt{2}}{3}\,M_\psi\gamma\omega_0 q_0^2.
\label{eq:friction_exact}
\end{align}

Using $\gamma = \omega_0/(2Q)$ and $\Delta V = \beta q_0^4/4 =
M_\psi\omega_0^2 q_0^2/2$ (matching the barrier to the mode energy), we
get:
\begin{equation}
W_{\mathrm{diss}} = \frac{4\sqrt{2}}{3}\cdot\frac{\omega_0}{2Q}
\cdot\frac{2\Delta V}{\omega_0} = \frac{4\sqrt{2}}{3}\cdot\frac{\Delta V}{Q}.
\label{eq:dissip_barrier}
\end{equation}

The numerical prefactor depends on the exact trajectory shape. For a
sinusoidal approximation $q(t) = q_0\sin(\omega_0 t)$, which is
appropriate for the harmonic oscillator limit near the well minimum:
\begin{equation}
\int_0^{\pi/\omega_0}\cos^2(\omega_0 t)\,dt = \frac{\pi}{2\omega_0},
\end{equation}
giving
\begin{equation}
W_{\mathrm{diss}} = 2M_\psi\gamma q_0^2\omega_0^2\cdot\frac{\pi}{2\omega_0}
= \frac{\pi M_\psi\gamma\omega_0 q_0^2}{1}
= \frac{\pi\Delta V}{Q}.
\label{eq:dissip_sinusoidal}
\end{equation}

We adopt the sinusoidal result~\eqref{eq:dissip_sinusoidal} as the
standard estimate, noting that the exact prefactor is $O(1)$ and
trajectory-dependent.

\textbf{Step 4: Substitution of reliability condition.}
For error rate below $P_{\mathrm{err}}$, the barrier must satisfy
$\Delta V \geq \alpha_{\mathrm{rel}}\kB T$ where $\alpha_{\mathrm{rel}}
= -\ln P_{\mathrm{err}}$. For $P_{\mathrm{err}} = 10^{-15}$:
$\alpha_{\mathrm{rel}} \approx 34.5$; for $P_{\mathrm{err}} =
10^{-22}$: $\alpha_{\mathrm{rel}} \approx 50.7$. Using the conservative
value $\alpha_{\mathrm{rel}} = 40$:
\begin{equation}
\langle W_{\mathrm{diss}}\rangle = \frac{\pi\alpha_{\mathrm{rel}}\kB T}{Q}.
\end{equation}

The entropy production is $\Delta S_{\mathrm{op}} = \langle
W_{\mathrm{diss}}\rangle/T = \pi\alpha_{\mathrm{rel}}\kB/Q$.
\end{proof}

\subsection{The Sub-Landauer Condition}\label{sec:sub_landauer_condition}

\begin{corollary}[Sub-Landauer Threshold]
\label{cor:sub_landauer}
A reversible $\psi$-field logic gate dissipates less than the Landauer
bound when
\begin{equation}
\boxed{
Q > Q_{\mathrm{crit}} \equiv \frac{\pi\alpha_{\mathrm{rel}}}{\ln 2}
}
\label{eq:Q_critical}
\end{equation}

For $\alpha_{\mathrm{rel}} = 40$: $Q_{\mathrm{crit}} \approx 181.3$.
For $\alpha_{\mathrm{rel}} = 50$: $Q_{\mathrm{crit}} \approx 226.6$.
\end{corollary}

\begin{proof}
Set $\langle W_{\mathrm{diss}}\rangle < \kB T\ln 2$ and solve for $Q$:
\begin{equation}
\frac{\pi\alpha_{\mathrm{rel}}\kB T}{Q} < \kB T\ln 2
\quad\Rightarrow\quad
Q > \frac{\pi\alpha_{\mathrm{rel}}}{\ln 2}.
\end{equation}
\end{proof}

\begin{remark}[No Violation of the Second Law]
This result does \textbf{not} violate the second law of thermodynamics
or Landauer's principle. The Landauer bound applies to logically
\emph{irreversible} operations (information erasure). Reversible gates
(NOT, SWAP, Toffoli, Fredkin) are bijective maps that erase no
information. The Jarzynski equality~\eqref{eq:jarzynski} with $\Delta F
= 0$ (Lemma~\ref{lem:zero_dF}) gives $\langle W\rangle \geq 0$, which
is satisfied. The dissipation $\pi\alpha_{\mathrm{rel}}\kB T/Q > 0$ is
positive, consistent with the second law. The point is simply that this
dissipation can be made arbitrarily small by increasing $Q$, because no
thermodynamic lower bound exists for reversible operations.
\end{remark}

\subsection{Entropy Production as a Function of $Q$}\label{sec:entropy_vs_Q}

The dissipation ratio relative to the Landauer bound is:
\begin{equation}
\mathcal{R}(Q) \equiv \frac{\langle W_{\mathrm{diss}}\rangle}{\kB T\ln 2}
= \frac{\pi\alpha_{\mathrm{rel}}}{Q\ln 2}.
\label{eq:dissipation_ratio}
\end{equation}

Table~\ref{tab:dissipation_vs_Q} gives numerical values for
representative quality factors.

\begin{table}[h]
\centering
\caption{Dissipation per reversible operation vs.\ cavity $Q$
($\alpha_{\mathrm{rel}} = 40$, $T = 1.5\,\mathrm{K}$).}
\label{tab:dissipation_vs_Q}
\begin{tabular}{@{}lccl@{}}
\toprule
$Q$ & $\mathcal{R}(Q)$ & $E_{\mathrm{diss}}$ (J) & Status \\
\midrule
$10^1$ & $18.1$ & $2.6\times 10^{-22}$ & Above Landauer \\
$10^2$ & $1.81$ & $2.6\times 10^{-23}$ & Above Landauer \\
$182$ & $1.00$ & $1.4\times 10^{-23}$ & At Landauer \\
$10^3$ & $0.181$ & $2.6\times 10^{-24}$ & Sub-Landauer \\
$10^5$ & $1.81\times 10^{-3}$ & $2.6\times 10^{-26}$ & $550\times$ below \\
$10^8$ & $1.81\times 10^{-6}$ & $2.6\times 10^{-29}$ & $5.5\times 10^5\times$ below \\
$10^{10}$ & $1.81\times 10^{-8}$ & $2.6\times 10^{-31}$ & $5.5\times 10^7\times$ below \\
$10^{11}$ & $1.81\times 10^{-9}$ & $2.6\times 10^{-32}$ & State-of-art SRF \\
\bottomrule
\end{tabular}
\end{table}

\subsection{Consistency Check via the Crooks Theorem}\label{sec:crooks_check}

As an independent verification, we apply the Crooks fluctuation
theorem~\eqref{eq:crooks}. For a reversible gate with $\Delta F = 0$,
the theorem reduces to:
\begin{equation}
\frac{P_F(W)}{P_R(-W)} = e^{W/\kB T}.
\label{eq:crooks_rev}
\end{equation}

This means the work distribution is symmetric about $W = 0$ in the
logarithmic ratio sense. The probability of observing negative work
(spontaneous reverse operation) is:
\begin{equation}
P(W < 0) = \int_{-\infty}^0 P_F(W)\,dW.
\end{equation}

For a Gaussian work distribution $P_F(W) = (2\pi\sigma_W^2)^{-1/2}
\exp[-(W - \langle W\rangle)^2/(2\sigma_W^2)]$ with $\langle W\rangle
= \pi\alpha_{\mathrm{rel}}\kB T/Q$ and $\sigma_W^2 = 2\kB T\langle
W\rangle$ (from the fluctuation-dissipation relation), the probability
of negative work is:
\begin{equation}
P(W < 0) = \frac{1}{2}\,\mathrm{erfc}\!\left(\sqrt{\frac{\langle W\rangle}{4\kB T}}\right)
= \frac{1}{2}\,\mathrm{erfc}\!\left(\sqrt{\frac{\pi\alpha_{\mathrm{rel}}}{4Q}}\right).
\end{equation}

For $Q = 10^{10}$ and $\alpha_{\mathrm{rel}} = 40$: $P(W < 0) \approx
0.4999944$, meaning the work distribution is almost perfectly symmetric
about zero---as expected for a nearly reversible process. The second law
is satisfied in the mean ($\langle W\rangle > 0$) even though individual
trajectories frequently have $W < 0$.

\subsection{Extension to Irreversible Gates}\label{sec:irreversible}

For a logically irreversible gate (AND, OR), the operation erases one
bit of input information. The free energy change is now:
\begin{equation}
\Delta F_{\mathrm{irrev}} = -\kB T\ln 2
\end{equation}
(the free energy decreases because the phase space contracts).

The Jarzynski equality gives:
\begin{equation}
\langle W\rangle \geq \Delta F = -\kB T\ln 2,
\end{equation}
which means $\langle W\rangle + \kB T\ln 2 \geq 0$, i.e., the total
energy cost (work plus heat absorbed from environment) is at least
$\kB T\ln 2$. The total dissipation is:
\begin{equation}
E_{\mathrm{irrev}} = \kB T\ln 2 + \frac{\pi\alpha_{\mathrm{rel}}\kB T}{Q},
\label{eq:irrev_dissipation}
\end{equation}
where the first term is the irreducible Landauer cost and the second is
the friction cost. For high-$Q$ cavities, the Landauer term dominates.

This motivates implementing all computation using reversible Toffoli or
Fredkin gates with ancilla bits, where the only entropy production comes
from the friction term and can be made arbitrarily small.

%=============================================================================
\section{Quantitative Literature Comparison}\label{sec:literature}
%=============================================================================

We compare $\psi$-field computing against 12 experimental and
theoretical benchmarks from the literature.

\subsection{Experimental Tests of Landauer's Principle}

\subsubsection{B\'er\'ut et al.\ (2012): Colloidal Particle Erasure}

B\'er\'ut et al.~\cite{Berut2012} performed the first direct
experimental test of Landauer's principle using a colloidal silica
particle ($d = 2\,\mu\mathrm{m}$) confined in a double-well optical
potential. By measuring the heat dissipated during one-bit erasure
protocols of varying duration $\tau$, they confirmed that the mean
dissipated heat satisfies $\langle Q_{\mathrm{diss}}\rangle \geq \kB T
\ln 2$ and approaches the bound for quasi-static protocols ($\tau \to
\infty$).

\paragraph{Comparison to $\psi$-field.}
The B\'er\'ut experiment operates at $T = 300\,\mathrm{K}$ with
effective $Q \sim 10$ (highly overdamped colloidal dynamics, viscous
friction dominates). The protocol duration was $\tau \sim 1$--$25\,\mathrm{s}$
for erasure at $\sim \kB T\ln 2$. A $\psi$-field cavity at $Q \sim
10^{10}$ and $T = 1.5\,\mathrm{K}$ achieves:
\begin{equation}
\frac{E_{\psi\text{-field}}}{E_{\text{B\'er\'ut}}} \sim \frac{1.5 \times 10^{-10}}{300 \times 10^{-0}} \sim 10^{-12},
\end{equation}
i.e., $10^{12}$ times less dissipation per operation (reversible
operation, bypassing Landauer entirely), at $10^9$ times faster speed.

\subsubsection{Jun et al.\ (2014): Feedback-Controlled Erasure}

Jun, Gavrilov, and Bhcek~\cite{Jun2014} used feedback-controlled
erasure of a colloidal particle in an asymmetric double well,
demonstrating that the Landauer bound can be approached from below for
\emph{partial} erasure (reducing information content by less than one
bit). They measured dissipation down to $\sim 0.7\,\kB T\ln 2$ for
partial erasure, confirming the generalized Landauer bound
$Q_{\mathrm{diss}} \geq \kB T\,\Delta I$ where $\Delta I$ is the mutual
information reduction.

\paragraph{Comparison.} This experiment confirms that the Landauer bound
is tight for irreversible operations. The $\psi$-field approach bypasses
this constraint entirely for reversible operations; for irreversible
operations, it saturates the bound with negligible excess.

\subsubsection{Hong et al.\ (2016): Nanomagnetic Erasure}

Hong, Lambson, Dhuey, and Bokor~\cite{Hong2016} tested Landauer's
principle in a nanomagnetic system, erasing single-domain
nanomagnets. They observed dissipation approaching $\kB T\ln 2$ for
slow erasure, with excess dissipation following a $1/\tau$ scaling
consistent with linear response theory.

\paragraph{Comparison.} The nanomagnetic system operates at $Q_{\mathrm{eff}}
\sim 100$. The $\psi$-field platform exceeds this by $10^8$ in $Q$,
translating to $10^8\times$ lower friction losses.

\subsubsection{Gaudenzi et al.\ (2018): Molecular Erasure}

Gaudenzi et al.~\cite{Gaudenzi2018} demonstrated Landauer erasure at
the single-molecule level using a molecular junction. They measured
dissipation consistent with the Landauer bound at $T = 0.5$--$10\,
\mathrm{K}$, accessing the quantum regime where $\hbar\omega \sim \kB T$.

\paragraph{Comparison.} This is the closest experimental system to
$\psi$-field computing in temperature regime. However, the molecular
junction has $Q \sim 10^2$--$10^3$, far below SRF cavities.

\subsection{Reversible Computing Proposals}

\subsubsection{Fredkin and Toffoli (1982): Conservative Logic}

Fredkin and Toffoli~\cite{FredkinToffoli1982} introduced conservative
logic, where every gate conserves the number of 1s in the bit string.
The Fredkin gate (controlled-SWAP) and Toffoli gate
(controlled-controlled-NOT) form universal sets for reversible computation.

\paragraph{$\psi$-field implementation.} The companion
paper~\cite{Alcock2025psiComputing} designs explicit $\psi$-field
Fredkin and Toffoli gates. The key advantage is that $\psi$-field
dynamics are \emph{naturally} Hamiltonian (energy-conserving), so
reversibility requires no special engineering---it is the default. The
energy overhead per gate is:
\begin{equation}
E_{\text{Toffoli}}^{(\psi)} = 3 \times \frac{\pi\alpha_{\mathrm{rel}}\kB T}{Q}
\approx 7.8\times 10^{-31}\,\mathrm{J}
\end{equation}
(three cavity modes involved, each contributing $\pi\Delta V/Q$).

\subsubsection{Bennett (1973, 1982): Reversible Turing Machines}

Bennett~\cite{Bennett1973} proved that any computation can be performed
reversibly with at most polynomial overhead in time and space. The key
technique is ``uncomputation'': after computing the output, reverse the
computation to erase intermediate (garbage) bits.

\paragraph{$\psi$-field advantage.} The Bennett overhead in electronic
reversible computing is dominated by the energy cost of signal
regeneration in lossy wires. In $\psi$-field computing, the signal
regeneration is \emph{free}---bistable cavity potentials automatically
restore degraded signals. The only overhead is the $3\times$ space
factor for the compute-copy-uncompute protocol.

\subsection{Adiabatic CMOS}

Adiabatic CMOS logic~\cite{Athas1994,Younis1994} uses slowly varying
clock signals to charge and discharge gate capacitances without
dissipating $CV^2/2$ per cycle. The dissipation scales as:
\begin{equation}
E_{\text{adiabatic}} \sim \frac{RC}{\tau}\cdot CV^2 = \frac{R C^2 V^2}{\tau},
\end{equation}
where $R$ is the resistance, $C$ is the capacitance, $V$ is the voltage
swing, and $\tau$ is the clock period. For $\tau \gg RC$, dissipation
approaches zero---but only in the limit of zero speed.

\paragraph{Comparison.} The adiabatic CMOS speed-dissipation tradeoff is
$E \propto 1/\tau$. The $\psi$-field tradeoff is $E = \pi\Delta V/Q$
independent of speed (the dissipation per cycle is fixed by $Q$, not by
how slowly the operation is performed). At $f = 1\,\mathrm{GHz}$:
\begin{equation}
\frac{E_{\text{adiabatic}}(1\,\mathrm{GHz})}{E_{\psi}(1\,\mathrm{GHz})} \sim \frac{10^{-18}}{3\times 10^{-31}} \sim 3\times 10^{12}.
\end{equation}

\subsection{Superconducting Logic}

\subsubsection{RSFQ/SFQ Logic}

Rapid Single Flux Quantum (RSFQ) logic~\cite{Likharev1991} encodes
information as the presence or absence of a magnetic flux quantum
$\Phi_0 = h/(2e) \approx 2.07\times 10^{-15}\,\mathrm{Wb}$ in
superconducting loops with Josephson junctions. Each switching event
dissipates:
\begin{equation}
E_{\mathrm{SFQ}} = I_c\Phi_0 \approx (100\,\mu\mathrm{A})(2\times 10^{-15}\,\mathrm{Wb})
\approx 2\times 10^{-19}\,\mathrm{J},
\end{equation}
at speeds of $\sim 50$--$770\,\mathrm{GHz}$ demonstrated in circuits.

\paragraph{Comparison.} SFQ is intrinsically dissipative (the flux
quantum tunneling through a Josephson junction is an irreversible event).
The $\psi$-field advantage in energy per operation is:
\begin{equation}
\frac{E_{\mathrm{SFQ}}}{E_\psi} \sim \frac{2\times 10^{-19}}{3\times 10^{-31}} \sim 7\times 10^{11}.
\end{equation}
However, SFQ currently operates at $\sim 100\times$ higher speed.

\subsubsection{AQFP Logic}

Adiabatic Quantum Flux Parametron (AQFP) logic~\cite{Takeuchi2013,
Takeuchi2017} operates Josephson junctions in an adiabatic regime,
achieving energy per operation of:
\begin{equation}
E_{\mathrm{AQFP}} \approx 10^{-21}\text{--}10^{-20}\,\mathrm{J}
\end{equation}
at clock frequencies of $\sim 5$--$10\,\mathrm{GHz}$. This approaches
the Landauer limit at 4.2\,K ($\kB T\ln 2 \approx 4\times
10^{-23}\,\mathrm{J}$).

\paragraph{Comparison.} AQFP is currently the closest competitor to
$\psi$-field computing in energy efficiency. The comparison:
\begin{center}
\begin{tabular}{lcc}
\toprule
Metric & AQFP & $\psi$-field \\
\midrule
$E$/op (J) & $10^{-21}$ & $3\times 10^{-31}$ \\
Speed (GHz) & 5--10 & 1--10 \\
$T$ (K) & 4.2 & 1.5 \\
Reversible? & Adiabatic & Hamiltonian \\
\bottomrule
\end{tabular}
\end{center}

The $\psi$-field advantage is $\sim 10^{10}\times$ in energy per
operation, though AQFP has far more mature fabrication.

\subsection{Photonic Computing}

\subsubsection{Photonic Tensor Cores}

Silicon photonic mesh networks~\cite{Shen2017,Harris2018} implement
matrix-vector multiplication using Mach-Zehnder interferometer (MZI)
arrays. The energy per multiply-accumulate (MAC) operation is:
\begin{equation}
E_{\mathrm{MZI}} \sim 10^{-15}\text{--}10^{-17}\,\mathrm{J/MAC}
\end{equation}
including photodetection and DAC/ADC conversion, at speeds $\sim
10\,\mathrm{GHz}$.

\paragraph{Comparison.} Photonic tensor cores face three fundamental
limitations that $\psi$-field computing avoids: (i)~optical-electronic
conversion overhead at each stage; (ii)~weak optical nonlinearities
($\chi^{(3)} \sim 10^{-22}\,\mathrm{m}^2/\mathrm{V}^2$) limiting
cascadability; (iii)~progressive signal degradation without
regeneration. The $\psi$-field bistable restoration provides free signal
regeneration.

\subsubsection{Photonic Reservoir Computing}

Larger et al.~\cite{Larger2012} demonstrated photonic reservoir
computing using a single nonlinear node with time-delay feedback,
achieving NARMA-10 benchmark performance comparable to digital
implementations. The energy efficiency was limited by the
optoelectronic transduction.

\paragraph{Comparison.} The $\psi$-field reservoir (a multi-mode cavity
with nonlinear $W(y)$ interaction) provides richer dynamics than a
single-node time-delay loop, potentially achieving higher computational
capacity per physical node.

\subsection{Neuromorphic Computing}

\subsubsection{Intel Loihi 2}

The Loihi 2 neuromorphic processor~\cite{Davies2018,Orchard2021}
implements 128 neuromorphic cores, each with 1024 spiking neurons, at
$\sim 15\,\mathrm{pJ}$ per synaptic operation. The total chip power is
$\sim 1\,\mathrm{W}$ at full utilization.

\paragraph{Quantitative comparison.}
\begin{equation}
\frac{E_{\mathrm{Loihi}}}{E_\psi} = \frac{1.5\times 10^{-11}}{3\times 10^{-31}} = 5\times 10^{19}.
\end{equation}

Even accounting for cryogenic overhead (COP $\sim 10^{-3}$ at 1.5\,K),
the $\psi$-field system-level efficiency advantage is $\sim 5\times
10^{16}$ per synaptic operation.

\subsubsection{IBM TrueNorth}

TrueNorth~\cite{Merolla2014} implements 4096 neurosynaptic cores with
$10^6$ neurons and $2.56\times 10^8$ synapses at 65\,mW total power
($\sim 26\,\mathrm{pJ}/\mathrm{synaptic\,event}$). It operates at
1\,kHz biological-time clock.

\subsubsection{Memristive Crossbar Arrays}

Memristive crossbar arrays~\cite{Prezioso2015,Li2018Nature} implement
matrix-vector multiplication in $O(1)$ time using Ohm's law and
Kirchhoff's current summation. Energy per MAC is $\sim 10^{-14}$--$10^{-12}\,\mathrm{J}$,
limited by the read current and wire resistance.

\paragraph{Comparison.} Memristive arrays share the analog
matrix-vector multiplication advantage of $\psi$-field computing but are
fundamentally dissipative (resistive current flow). The $\psi$-field
advantage is $10^{17}$--$10^{19}\times$ in energy per MAC.

\subsection{Summary Comparison Table}

\begin{table*}[t]
\centering
\caption{Comprehensive comparison of $\psi$-field computing against 12
benchmark technologies. Energy per operation, operating speed, operating
temperature, and key advantages/limitations.}
\label{tab:comprehensive}
\begin{tabular}{@{}lccccl@{}}
\toprule
\textbf{Technology} & \textbf{$E$/op (J)} & \textbf{Speed} & \textbf{$T$ (K)} &
\textbf{$E/\EL(T)$} & \textbf{Key limitation} \\
\midrule
7\,nm CMOS & $\sim 10^{-17}$ & 5\,GHz & 300 & $3500$ & Joule heating \\
2\,nm CMOS (proj.) & $\sim 10^{-18}$ & 5\,GHz & 300 & $350$ & Joule heating \\
Adiabatic CMOS & $\sim 10^{-18}$ & 100\,MHz & 300 & $350$ & Speed-energy tradeoff \\
RSFQ/SFQ & $\sim 2\times 10^{-19}$ & 100\,GHz & 4 & $5000$ & Intrinsic dissipation \\
AQFP & $\sim 10^{-21}$ & 5\,GHz & 4.2 & $25$ & Clock distribution \\
Photonic (MZI) & $\sim 10^{-17}$ & 10\,GHz & 300 & $3500$ & Weak nonlinearity \\
Photonic reservoir & $\sim 10^{-14}$ & 1\,GHz & 300 & $10^6$ & OE conversion \\
Intel Loihi 2 & $\sim 1.5\times 10^{-11}$ & 1\,GHz & 300 & $5\times 10^9$ & Digital overhead \\
IBM TrueNorth & $\sim 2.6\times 10^{-11}$ & 1\,kHz & 300 & $9\times 10^9$ & Low clock rate \\
Memristive array & $\sim 10^{-14}$ & 100\,MHz & 300 & $3\times 10^6$ & Device variability \\
Colloidal (B\'er\'ut) & $\sim 3\times 10^{-21}$ & 0.1\,Hz & 300 & $1$ & Macroscopic, slow \\
\midrule
$\psi$-field (reversible) & $\sim 3\times 10^{-31}$ & 1\,GHz & 1.5 & $2\times 10^{-8}$ & Cryogenic, size \\
\bottomrule
\end{tabular}
\end{table*}

%=============================================================================
\section{Complete 4-Bit ALU Design}\label{sec:alu}
%=============================================================================

We design a complete 4-bit arithmetic logic unit (ALU) using $\psi$-field
reversible logic gates. The ALU implements: addition, subtraction, AND,
OR, XOR, NOT, and left/right shift.

\subsection{Gate Primitives}\label{sec:alu_primitives}

All ALU operations are built from the following reversible $\psi$-field
gate primitives:

\paragraph{Toffoli gate (CCNOT).} Three cavities: controls $A, B$,
target $C$. Operation: $C \to C \oplus (A \wedge B)$.
\begin{equation}
H_{\mathrm{Tof}} = H_A + H_B + H_C + \Lambda_4 q_A q_B(q_C^2 - q_0^2) + h_T q_A q_B q_C.
\end{equation}
Energy: $E_{\mathrm{Tof}} = 3\pi\alpha_{\mathrm{rel}}\kB T/Q$.
Latency: $\tau_{\mathrm{Tof}} = 3\pi/\omega_0$ (three sequential mode
interactions).

\paragraph{CNOT gate.} Two cavities: control $A$, target $B$. Operation:
$B \to B \oplus A$. Implemented as Toffoli with one control fixed at
$|1\rangle$.
Energy: $E_{\mathrm{CNOT}} = 2\pi\alpha_{\mathrm{rel}}\kB T/Q$.
Latency: $\tau_{\mathrm{CNOT}} = 2\pi/\omega_0$.

\paragraph{NOT gate.} Single cavity with inverting bias pulse.
Energy: $E_{\mathrm{NOT}} = \pi\alpha_{\mathrm{rel}}\kB T/Q$.
Latency: $\tau_{\mathrm{NOT}} = \pi/\omega_0$.

\paragraph{Fredkin gate (CSWAP).} Control $A$, targets $B, C$. Swaps
$B \leftrightarrow C$ when $A = |1\rangle$.
Energy: $E_{\mathrm{Fred}} = 3\pi\alpha_{\mathrm{rel}}\kB T/Q$.

\subsection{4-Bit Ripple-Carry Adder}\label{sec:adder}

A 4-bit reversible adder uses the Cuccaro et
al.~\cite{Cuccaro2004} construction, requiring:
\begin{itemize}
\item Inputs: $A = a_3 a_2 a_1 a_0$, $B = b_3 b_2 b_1 b_0$, carry-in $c_0$.
\item Outputs: $S = s_3 s_2 s_1 s_0$, carry-out $c_4$.
\item Ancilla: 1 bit.
\end{itemize}

\paragraph{Gate count.} The Cuccaro adder for $n$ bits requires $2n - 1$
Toffoli gates and $2n$ CNOT gates. For $n = 4$:
\begin{itemize}
\item 7 Toffoli gates
\item 8 CNOT gates
\item Total: 15 reversible gates
\end{itemize}

\paragraph{Energy per addition.}
\begin{align}
E_{\mathrm{add}} &= 7\,E_{\mathrm{Tof}} + 8\,E_{\mathrm{CNOT}} \nonumber\\
&= (7\times 3 + 8\times 2)\frac{\pi\alpha_{\mathrm{rel}}\kB T}{Q} \nonumber\\
&= \frac{37\pi\alpha_{\mathrm{rel}}\kB T}{Q}.
\label{eq:E_add}
\end{align}

For $Q = 10^{10}$, $T = 1.5\,\mathrm{K}$, $\alpha_{\mathrm{rel}} = 40$:
\begin{equation}
E_{\mathrm{add}} = \frac{37\pi \times 40 \times 1.38\times 10^{-23} \times 1.5}{10^{10}}
\approx 9.6\times 10^{-30}\,\mathrm{J}.
\end{equation}

\paragraph{Latency.}
The critical path through a 4-bit ripple-carry adder has depth 13
(sequential Toffoli and CNOT operations along the carry chain):
\begin{equation}
\tau_{\mathrm{add}} = 13 \times \frac{\pi}{\omega_0}
= \frac{13}{2f_0} \approx 6.5\,\mathrm{ns} \quad (f_0 = 1\,\mathrm{GHz}).
\end{equation}

\paragraph{Throughput (pipelined).}
With pipeline registers between stages:
\begin{equation}
f_{\mathrm{throughput}} = f_0 = 1\,\mathrm{GHz} = 10^9\,\mathrm{additions/s}.
\end{equation}

\subsection{4-Bit Subtraction}\label{sec:subtractor}

Subtraction $A - B$ is implemented as addition of the two's complement:
$A + \overline{B} + 1$. This requires:
\begin{itemize}
\item 4 NOT gates (to compute $\overline{B}$)
\item 1 adder (15 gates, as above) with carry-in $c_0 = 1$
\item Total: 19 reversible gates
\end{itemize}

\begin{equation}
E_{\mathrm{sub}} = 4\,E_{\mathrm{NOT}} + E_{\mathrm{add}}
= \left(4 + 37\right)\frac{\pi\alpha_{\mathrm{rel}}\kB T}{Q}
= \frac{41\pi\alpha_{\mathrm{rel}}\kB T}{Q}
\approx 1.1\times 10^{-29}\,\mathrm{J}.
\end{equation}

\subsection{Bitwise Logic Operations}\label{sec:bitwise}

\paragraph{4-bit AND.}
Reversible AND uses a Toffoli gate per bit with a fresh ancilla:
$(a_i, b_i, 0) \to (a_i, b_i, a_i \wedge b_i)$. For 4 bits:
\begin{equation}
E_{\mathrm{AND}} = 4\,E_{\mathrm{Tof}} = \frac{12\pi\alpha_{\mathrm{rel}}\kB T}{Q}
\approx 3.1\times 10^{-30}\,\mathrm{J}.
\end{equation}
Latency: $\tau_{\mathrm{AND}} = 4\tau_{\mathrm{Tof}} = 12\pi/\omega_0
\approx 6\,\mathrm{ns}$ (parallel execution possible, reducing to
$\tau_{\mathrm{Tof}} = 1.5\,\mathrm{ns}$).

\paragraph{4-bit OR.}
Using De Morgan: $a \vee b = \overline{\overline{a} \wedge
\overline{b}}$. Requires 3 NOT gates and 1 Toffoli per bit (with
uncomputation):
\begin{equation}
E_{\mathrm{OR}} = 4(3\,E_{\mathrm{NOT}} + E_{\mathrm{Tof}})
= 4\left(\frac{3 + 3}{1}\right)\frac{\pi\alpha_{\mathrm{rel}}\kB T}{Q}
= \frac{24\pi\alpha_{\mathrm{rel}}\kB T}{Q}
\approx 6.2\times 10^{-30}\,\mathrm{J}.
\end{equation}

\paragraph{4-bit XOR.}
XOR is natively reversible: $a \oplus b$ is computed by a single CNOT
per bit: $(a_i, b_i) \to (a_i, a_i \oplus b_i)$.
\begin{equation}
E_{\mathrm{XOR}} = 4\,E_{\mathrm{CNOT}} = \frac{8\pi\alpha_{\mathrm{rel}}\kB T}{Q}
\approx 2.1\times 10^{-30}\,\mathrm{J}.
\end{equation}
Latency: $\tau_{\mathrm{XOR}} = \tau_{\mathrm{CNOT}} = 1\,\mathrm{ns}$
(all 4 CNOTs execute in parallel).

\paragraph{4-bit NOT.}
Four parallel NOT gates:
\begin{equation}
E_{\mathrm{NOT4}} = 4\,E_{\mathrm{NOT}} = \frac{4\pi\alpha_{\mathrm{rel}}\kB T}{Q}
\approx 1.0\times 10^{-30}\,\mathrm{J}.
\end{equation}
Latency: $\pi/\omega_0 = 0.5\,\mathrm{ns}$.

\paragraph{4-bit left shift.}
A reversible left shift by 1 position is a cyclic permutation:
$(a_3, a_2, a_1, a_0) \to (a_2, a_1, a_0, 0)$ with carry-out $a_3$.
Implemented using 3 SWAP operations (each SWAP = 3 CNOTs):
\begin{equation}
E_{\mathrm{shift}} = 3 \times 3\,E_{\mathrm{CNOT}} = \frac{18\pi\alpha_{\mathrm{rel}}\kB T}{Q}
\approx 4.7\times 10^{-30}\,\mathrm{J}.
\end{equation}

\subsection{ALU Summary and Comparison}\label{sec:alu_comparison}

\begin{table}[h]
\centering
\caption{4-bit ALU operation comparison: $\psi$-field (SRF, $Q =
10^{10}$, $T = 1.5\,\mathrm{K}$) vs.\ 7\,nm CMOS, AQFP, and photonic.}
\label{tab:alu_comparison}
\begin{tabular}{@{}lccccc@{}}
\toprule
 & \multicolumn{2}{c}{$\psi$-field} & \multicolumn{1}{c}{7\,nm} & \multicolumn{1}{c}{AQFP} & \multicolumn{1}{c}{Photonic} \\
\textbf{Op} & $E$ (J) & $\tau$ (ns) & $E$ (J) & $E$ (J) & $E$ (J) \\
\midrule
ADD & $10^{-29}$ & 6.5 & $10^{-16}$ & $10^{-20}$ & $10^{-16}$ \\
SUB & $10^{-29}$ & 7.5 & $10^{-16}$ & $10^{-20}$ & $10^{-16}$ \\
AND & $3\times 10^{-30}$ & 1.5 & $10^{-17}$ & $10^{-21}$ & $10^{-17}$ \\
OR & $6\times 10^{-30}$ & 3 & $10^{-17}$ & $10^{-21}$ & $10^{-17}$ \\
XOR & $2\times 10^{-30}$ & 1 & $10^{-17}$ & $10^{-21}$ & $10^{-17}$ \\
NOT & $10^{-30}$ & 0.5 & $10^{-17}$ & $10^{-21}$ & $10^{-17}$ \\
SHIFT & $5\times 10^{-30}$ & 3 & $10^{-17}$ & $10^{-21}$ & --- \\
\bottomrule
\end{tabular}
\end{table}

\paragraph{Cavity count.}
The complete 4-bit ALU (with multiplexing between operations) requires:
\begin{itemize}
\item 8 input cavities (4-bit $A$, 4-bit $B$)
\item 5 output cavities (4-bit result + carry/flags)
\item 19 intermediate/ancilla cavities (for reversible computation)
\item 15 coupling/nonlinear junction elements
\item \textbf{Total: 47 cavities}
\end{itemize}

At 1.3\,GHz (23\,cm per cavity), the physical footprint is approximately
$1\,\mathrm{m} \times 0.5\,\mathrm{m} \times 0.5\,\mathrm{m}$. At
10\,GHz (1.5\,cm per cavity), this shrinks to $\sim 15\,\mathrm{cm}
\times 8\,\mathrm{cm} \times 8\,\mathrm{cm}$.

\subsection{4-Bit ALU Circuit Architecture}\label{sec:alu_circuit}

The ALU is organized as follows:

\begin{center}
\begin{tikzpicture}[scale=0.55, every node/.style={font=\scriptsize}]
  % Input registers
  \draw[thick, fill=blue!10] (-5,3) rectangle (-3,5);
  \node at (-4,4) {$A_{3:0}$};
  \node[above] at (-4,5.2) {Input A};

  \draw[thick, fill=blue!10] (-5,-1) rectangle (-3,1);
  \node at (-4,0) {$B_{3:0}$};
  \node[below] at (-4,-1.2) {Input B};

  % Op-code
  \draw[thick, fill=yellow!20] (-5,1.5) rectangle (-3,2.5);
  \node at (-4,2) {OpCode};

  % Routing network
  \draw[thick, fill=orange!10] (-1,0) rectangle (3,4);
  \node at (1,3.5) {\textbf{Fredkin}};
  \node at (1,2.5) {\textbf{Routing}};
  \node at (1,1.5) {\textbf{Network}};
  \node at (1,0.5) {(op select)};

  % Functional units
  \draw[thick, fill=green!10] (5,4.5) rectangle (9,6);
  \node at (7,5.25) {Adder/Sub};

  \draw[thick, fill=green!10] (5,2.5) rectangle (9,4);
  \node at (7,3.25) {AND/OR/XOR};

  \draw[thick, fill=green!10] (5,0.5) rectangle (9,2);
  \node at (7,1.25) {Shift/NOT};

  % Output
  \draw[thick, fill=red!10] (11,1.5) rectangle (14,4.5);
  \node at (12.5,3) {Output};
  \node at (12.5,2) {$S_{3:0}$};

  % Arrows
  \draw[-Stealth,thick] (-3,4) -- (-1,3);
  \draw[-Stealth,thick] (-3,0) -- (-1,1);
  \draw[-Stealth,thick] (-3,2) -- (-1,2);
  \draw[-Stealth,thick] (3,3.5) -- (5,5.25);
  \draw[-Stealth,thick] (3,2) -- (5,3.25);
  \draw[-Stealth,thick] (3,0.5) -- (5,1.25);
  \draw[-Stealth,thick] (9,5.25) -- (11,3.5);
  \draw[-Stealth,thick] (9,3.25) -- (11,3);
  \draw[-Stealth,thick] (9,1.25) -- (11,2.5);
\end{tikzpicture}
\end{center}

The operation selection is performed by a Fredkin routing network
controlled by a 3-bit opcode. This is itself a reversible operation,
adding no irreversible dissipation. The opcode cavities steer input data
to the appropriate functional unit via controlled-SWAP operations.

%=============================================================================
\section{Radiation Hardness Analysis}\label{sec:radiation}
%=============================================================================

\subsection{Single-Event Upset Physics}\label{sec:seu_physics}

A single-event upset (SEU) occurs when an energetic particle (proton,
heavy ion, or secondary from nuclear interaction) deposits sufficient
energy in a memory cell to change its logical state. In semiconductor
SRAM, the critical charge is:
\begin{equation}
Q_{\mathrm{crit}}^{\mathrm{SRAM}} = C_{\mathrm{node}} \times V_{\mathrm{DD}}
\sim (1\,\mathrm{fF})(0.7\,\mathrm{V}) = 0.7\,\mathrm{fC} \approx 4400\,e^-.
\label{eq:Qcrit_SRAM}
\end{equation}

A particle traversing the sensitive volume deposits charge via
ionization:
\begin{equation}
\Delta Q = \frac{\mathrm{LET} \times \rho \times d}{W}
\end{equation}
where LET is the linear energy transfer ($\mathrm{MeV\cdot
cm}^2/\mathrm{mg}$), $\rho$ is the density, $d$ is the sensitive depth,
and $W \approx 3.6\,\mathrm{eV}$ is the electron-hole pair creation
energy in silicon.

The SEU cross-section is:
\begin{equation}
\sigma_{\mathrm{SEU}}^{\mathrm{SRAM}} = A_{\mathrm{sensitive}} \times
\left[1 - \exp\!\left(-\frac{\mathrm{LET}}{\mathrm{LET}_{\mathrm{th}}}\right)\right]^s,
\end{equation}
where $A_{\mathrm{sensitive}}$ is the sensitive area and
$\mathrm{LET}_{\mathrm{th}}$ is the threshold LET. For modern SRAM at
7\,nm:
\begin{equation}
\sigma_{\mathrm{SEU}}^{\mathrm{SRAM}} \approx 10^{-14}\text{--}10^{-12}\,\mathrm{cm}^2/\mathrm{bit}.
\end{equation}

\subsection{$\psi$-Field Memory: Upset Mechanism}\label{sec:psi_seu}

For a $\psi$-field memory cell, the state is encoded in the collective
mode amplitude $q$ of a macroscopic cavity. An incident particle
interacts with the cavity through two channels:

\paragraph{(a) Direct EM coupling.}
A charged particle traversing the cavity generates a transient
electromagnetic field that couples to the cavity modes. The deposited
energy per mode is:
\begin{equation}
\delta E_n = \frac{(Ze)^2\omega_n}{4\epsilon_0 V_{\mathrm{cav}}}
\left|\int_{\mathrm{path}} \phi_n(\mathbf{x})\,ds\right|^2,
\label{eq:em_coupling}
\end{equation}
where $Ze$ is the particle charge, $\omega_n$ is the mode frequency,
$V_{\mathrm{cav}}$ is the cavity volume, and $\phi_n$ is the mode
function.

\paragraph{(b) Nuclear recoil.}
Nuclear scattering in the cavity walls creates lattice damage and
secondary particles, but these do not directly couple to the
electromagnetic mode amplitude.

\subsection{Mode-Number Suppression}\label{sec:mode_suppression}

The critical insight is that a $\psi$-field memory bit stored in
$N_{\mathrm{mode}}$ coherently excited modes requires \emph{all} modes
to be simultaneously perturbed beyond their individual thresholds for a
bit flip to occur.

\begin{theorem}[Exponential SEU Suppression]
\label{thm:seu_suppression}
The SEU cross-section for a $\psi$-field memory cell using
$N_{\mathrm{mode}}$ participating cavity modes is
\begin{equation}
\boxed{
\sigma_{\mathrm{SEU}}^{(\psi)} = \sigma_{\mathrm{geom}}
\prod_{n=1}^{N_{\mathrm{mode}}} P_n(\mathrm{flip})
= \sigma_{\mathrm{geom}}\left(\frac{\delta E}{E_{\mathrm{barrier}}}\right)^{N_{\mathrm{mode}}}
}
\label{eq:seu_psi}
\end{equation}
where $\sigma_{\mathrm{geom}}$ is the geometric cross-section of the
cavity, $\delta E$ is the energy deposited per mode by the particle, and
$E_{\mathrm{barrier}} = \alpha_{\mathrm{rel}}\kB T$ is the barrier
height.
\end{theorem}

\begin{proof}
The mode amplitudes $\{q_n\}$ are coupled through the nonlinear
potential but fluctuate independently on short timescales (compared to
the coupling time). A bit flip requires transitioning the collective
state from one basin of attraction to the other, which requires
perturbing each participating mode by enough to overcome its contribution
to the barrier.

For a particle depositing total energy $\delta E_{\mathrm{tot}}$
distributed across modes, the fraction going to mode $n$ is:
\begin{equation}
\delta E_n = \delta E_{\mathrm{tot}} \times |\langle\phi_n|\delta\psi_{\mathrm{particle}}\rangle|^2.
\end{equation}

For a localized particle track (spatial extent $\ell_{\mathrm{track}}
\ll L_{\mathrm{cavity}}$), the overlap with mode $n$ scales as:
\begin{equation}
|\langle\phi_n|\delta\psi\rangle|^2 \sim \left(\frac{\ell_{\mathrm{track}}}{L}\right)^3
\times \frac{1}{N_{\mathrm{mode}}}
\end{equation}
because the mode functions are delocalized over the full cavity volume.

The probability that mode $n$ is flipped is:
\begin{equation}
P_n = \min\!\left(1,\;\frac{\delta E_n}{E_{\mathrm{barrier}}}\right).
\end{equation}

Since the mode flips are approximately independent events (different
spatial mode functions), the probability of a complete bit flip is:
\begin{equation}
P_{\mathrm{SEU}} = \prod_{n=1}^{N_{\mathrm{mode}}} P_n
\leq \left(\frac{\delta E_{\mathrm{tot}}}{N_{\mathrm{mode}}\,E_{\mathrm{barrier}}}\right)^{N_{\mathrm{mode}}},
\end{equation}
giving the exponential suppression.
\end{proof}

\subsection{Numerical Cross-Section Estimates}\label{sec:seu_numbers}

We compute $\sigma_{\mathrm{SEU}}^{(\psi)}$ for specific parameters.

\paragraph{Geometric cross-section.}
For a 1.3\,GHz cavity ($L = 23\,\mathrm{cm}$, $D = 18.7\,\mathrm{cm}$):
\begin{equation}
\sigma_{\mathrm{geom}} = \frac{\pi D^2}{4} \approx 275\,\mathrm{cm}^2.
\end{equation}

\paragraph{Energy deposition per mode.}
A 1\,GeV proton traversing the cavity deposits (via direct EM coupling
to the fundamental mode, Eq.~\eqref{eq:em_coupling}):
\begin{equation}
\delta E_1 \sim \frac{e^2 \omega_0}{4\epsilon_0 V_{\mathrm{cav}}}
\times L^2 \approx \frac{(1.6\times 10^{-19})^2 \times 8.2\times 10^9}
{4 \times 8.85\times 10^{-12} \times 0.006}
\approx 10^{-16}\,\mathrm{J}.
\end{equation}

\paragraph{Barrier height.}
At $T = 1.5\,\mathrm{K}$ with $\alpha_{\mathrm{rel}} = 50$:
\begin{equation}
E_{\mathrm{barrier}} = 50 \times 1.38\times 10^{-23} \times 1.5
= 1.04\times 10^{-21}\,\mathrm{J}.
\end{equation}

\paragraph{Ratio.}
\begin{equation}
\frac{\delta E_1}{E_{\mathrm{barrier}}} \sim \frac{10^{-16}}{10^{-21}} = 10^5.
\end{equation}

This means a single mode \emph{is} easily flipped by a GeV proton.
However, with $N_{\mathrm{mode}}$ modes, the energy per mode decreases:
for a localized track coupling primarily to one mode:
\begin{equation}
\frac{\delta E_n}{E_{\mathrm{barrier}}} \approx \frac{10^5}{N_{\mathrm{mode}}}.
\end{equation}

For $N_{\mathrm{mode}} > 10^5$, no individual mode is flipped. For
intermediate $N_{\mathrm{mode}}$:

\begin{table}[h]
\centering
\caption{$\psi$-field SEU cross-section vs.\ number of modes, compared
to SRAM and MRAM.}
\label{tab:seu}
\begin{tabular}{@{}lcc@{}}
\toprule
\textbf{Technology} & $N_{\mathrm{mode}}$ & $\sigma_{\mathrm{SEU}}$ (cm$^2$/bit) \\
\midrule
7\,nm SRAM & --- & $10^{-14}$--$10^{-12}$ \\
STT-MRAM & --- & $10^{-16}$--$10^{-14}$ \\
SOT-MRAM & --- & $10^{-18}$--$10^{-16}$ \\
\midrule
$\psi$-field, $N = 1$ & 1 & $275$ (trivially flipped) \\
$\psi$-field, $N = 3$ & 3 & $275 \times (3\times 10^{4})^{-3} \approx 10^{-11}$ \\
$\psi$-field, $N = 5$ & 5 & $275 \times (2\times 10^{4})^{-5} \approx 10^{-19}$ \\
$\psi$-field, $N = 10$ & 10 & $275 \times (10^{4})^{-10} \approx 10^{-38}$ \\
$\psi$-field, $N = 20$ & 20 & $< 10^{-75}$ \\
\bottomrule
\end{tabular}
\end{table}

\begin{remark}
With $N_{\mathrm{mode}} \geq 5$ participating modes, $\psi$-field
memory achieves SEU cross-sections below $10^{-19}\,\mathrm{cm}^2
/\mathrm{bit}$, surpassing the most radiation-hardened MRAM
technologies. With $N_{\mathrm{mode}} = 10$, the cross-section is
effectively zero for any physical particle flux. This makes $\psi$-field
memory inherently suited to space, nuclear, and high-energy physics
environments.
\end{remark}

\subsection{Design Guidelines for Radiation-Hard $\psi$-Memory}
\label{sec:rad_design}

To maximize radiation hardness:
\begin{enumerate}
\item \textbf{Multi-mode encoding}: Store each bit in $N_{\mathrm{mode}}
\geq 5$ coherently excited modes. The exponential suppression makes this
the dominant hardening mechanism.
\item \textbf{Distributed mode functions}: Use cavity geometries that
maximize the spatial delocalization of mode functions (e.g., confocal
cavities, whispering-gallery modes).
\item \textbf{Shielding}: Standard radiation shielding reduces the
incident flux. Combined with multi-mode encoding, even moderate shielding
($\sim 1\,\mathrm{mm}$ Al) provides extraordinary protection.
\item \textbf{Active scrubbing}: Periodic readout and refresh of stored
bits corrects any accumulated sub-threshold perturbations.
\end{enumerate}

%=============================================================================
\section{$\psi$-Field Neural Network Layer}\label{sec:neural}
%=============================================================================

\subsection{Layer Architecture}\label{sec:layer_arch}

We design a complete dense neural network layer mapping $N_{\mathrm{in}}$
input neurons to $N_{\mathrm{out}}$ output neurons.

\paragraph{Physical layout.}
\begin{itemize}
\item \textbf{Input layer}: $N_{\mathrm{in}}$ $\psi$-field cavities with
mode amplitudes $x_i \equiv q_i/q_0 \in [-1, 1]$.
\item \textbf{Coupling matrix}: $N_{\mathrm{in}} \times
N_{\mathrm{out}}$ coupling apertures implementing weights $w_{ji}
\equiv \kappa_{ji}/\kappa_0$ (normalized coupling strengths).
\item \textbf{Output layer}: $N_{\mathrm{out}}$ $\psi$-field cavities
with bistable activation producing outputs $y_j$.
\item \textbf{Bias}: External field gradients $b_j$ applied to output
cavities.
\end{itemize}

\subsection{Forward Pass}\label{sec:forward}

The forward pass computes:
\begin{equation}
\boxed{
y_j = \sigma\!\left(\sum_{i=1}^{N_{\mathrm{in}}} w_{ji}\,x_i + b_j\right),
\qquad j = 1, \ldots, N_{\mathrm{out}}
}
\label{eq:forward}
\end{equation}
where $\sigma(z) = \tanh(z/z_0)$ is the activation function arising
from the bistable cavity response (see below).

\paragraph{Physical mechanism.}
The input cavities have mode amplitudes $q_i = q_0 x_i$ representing
the input vector $\mathbf{x}$. Through the coupling apertures, the
$i$-th input cavity exerts a force on the $j$-th output cavity:
\begin{equation}
F_{j \leftarrow i} = -\frac{\partial J_{ji}}{\partial q_j}
= -\kappa_{ji}\,q_i = -\kappa_0 w_{ji}\,q_0\,x_i.
\end{equation}

The total input to output cavity $j$ is:
\begin{equation}
F_j^{\mathrm{total}} = \sum_i F_{j\leftarrow i} + F_j^{\mathrm{bias}}
= -\kappa_0 q_0\left(\sum_i w_{ji}\,x_i + b_j\right)
\equiv -\kappa_0 q_0\,z_j.
\label{eq:total_input}
\end{equation}

The output cavity $j$ has a double-well potential. Under the influence
of $F_j^{\mathrm{total}}$, it settles to a steady-state amplitude:
\begin{equation}
q_j^{\mathrm{out}} = q_0\,\tanh\!\left(\frac{\kappa_0 q_0\,z_j}{2\Delta V/q_0}\right)
= q_0\,\tanh\!\left(\frac{z_j}{z_0}\right),
\label{eq:tanh_activation}
\end{equation}
where $z_0 = 2\Delta V/(\kappa_0 q_0^2)$ is the activation
sensitivity scale, controlled by the barrier height and coupling
strength.

The output is $y_j = q_j^{\mathrm{out}}/q_0 = \tanh(z_j/z_0)$.

\paragraph{Computation time.}
The matrix-vector multiplication $\mathbf{z} = W\mathbf{x} + \mathbf{b}$
is performed by wave propagation from input to output cavities. The
computation time is:
\begin{equation}
\tau_{\mathrm{forward}} = \tau_{\mathrm{prop}} + \tau_{\mathrm{settle}}
= \frac{L_{\mathrm{coupling}}}{c_\psi} + \frac{\pi}{\omega_0}
\approx 1\text{--}2\,\mathrm{ns},
\end{equation}
\emph{independent of the layer size} $N_{\mathrm{in}}, N_{\mathrm{out}}$.
This is because all input-to-output couplings operate simultaneously---the
matrix multiplication is inherently parallel.

\paragraph{Energy per forward pass.}
Each output cavity undergoes one settling operation:
\begin{equation}
E_{\mathrm{forward}} = N_{\mathrm{out}} \times \frac{\pi\alpha_{\mathrm{rel}}\kB T}{Q}.
\end{equation}
For $N_{\mathrm{out}} = 1024$, $Q = 10^{10}$, $T = 1.5\,\mathrm{K}$:
\begin{equation}
E_{\mathrm{forward}} = 1024 \times 2.6\times 10^{-31}
\approx 2.7\times 10^{-28}\,\mathrm{J}.
\end{equation}

The number of MAC operations per forward pass is $N_{\mathrm{in}}
\times N_{\mathrm{out}}$. The energy per MAC is therefore:
\begin{equation}
E_{\mathrm{MAC}} = \frac{E_{\mathrm{forward}}}{N_{\mathrm{in}} \times N_{\mathrm{out}}}
= \frac{\pi\alpha_{\mathrm{rel}}\kB T}{Q\,N_{\mathrm{in}}}.
\end{equation}
For $N_{\mathrm{in}} = 1024$: $E_{\mathrm{MAC}} \approx 2.5\times
10^{-34}\,\mathrm{J}$.

\subsection{Backward Pass (Backpropagation)}\label{sec:backward}

Training requires computing gradients of a loss function $\mathcal{L}$
with respect to the weights $w_{ji}$. The backward pass computes:

\paragraph{Step 1: Output error.}
The error signal at the output layer is:
\begin{equation}
\delta_j = \frac{\partial\mathcal{L}}{\partial z_j}
= \frac{\partial\mathcal{L}}{\partial y_j}\cdot\sigma'(z_j/z_0)/z_0,
\label{eq:output_error}
\end{equation}
where $\sigma'(u) = 1 - \tanh^2(u) = \sech^2(u)$.

\paragraph{Physical implementation of $\delta_j$.}
The error signal $\delta_j$ is injected into the output cavities as an
external perturbation. The $\sech^2$ factor arises naturally from the
linearized response of the bistable cavity around its operating point:
\begin{equation}
\frac{\partial q_j^{\mathrm{out}}}{\partial F_j}
= \frac{q_0}{\kappa_0 q_0}\cdot\frac{1}{z_0}\,\sech^2\!\left(\frac{z_j}{z_0}\right)
= \frac{1}{\kappa_0 z_0}\,\sech^2\!\left(\frac{z_j}{z_0}\right).
\end{equation}

\paragraph{Step 2: Weight gradient.}
\begin{equation}
\boxed{
\frac{\partial\mathcal{L}}{\partial w_{ji}} = \delta_j\,x_i
}
\label{eq:weight_gradient}
\end{equation}

This is a \emph{local} computation: the gradient for weight $w_{ji}$
depends only on the error at output $j$ and the activation of input $i$.

\paragraph{Physical implementation.}
The weight gradient is computed by correlating the error signal $\delta_j$
in output cavity $j$ with the input signal $x_i$ in input cavity $i$.
This correlation is performed by a mediating cavity in the coupling
channel between $i$ and $j$:
\begin{equation}
\frac{d q_{\mathrm{med},ji}}{dt} = \eta\,\delta_j(t)\,x_i(t) - \lambda_d\,q_{\mathrm{med},ji},
\end{equation}
where the mediating cavity amplitude $q_{\mathrm{med},ji}$ accumulates
the correlation over the learning period. The decay term
$\lambda_d\,q_{\mathrm{med},ji}$ prevents runaway growth (weight
decay / L2 regularization).

\paragraph{Step 3: Error backpropagation.}
For a multi-layer network, the error must propagate backward through the
coupling matrix. The backpropagated error to input neuron $i$ is:
\begin{equation}
\boxed{
\delta_i^{(l)} = \sigma'\!\left(\frac{z_i^{(l)}}{z_0}\right)\cdot\frac{1}{z_0}
\sum_j w_{ji}\,\delta_j^{(l+1)}
}
\label{eq:backprop}
\end{equation}

\paragraph{Physical implementation.}
The backpropagated error $\delta_i^{(l)}$ is computed by injecting
$\delta_j^{(l+1)}$ into the output cavities and reading the resulting
perturbation at the input cavities through the \emph{same} coupling
apertures (exploiting the reciprocity of the coupling). The transposed
matrix-vector multiplication $\sum_j w_{ji}\delta_j$ is performed by
wave propagation in the reverse direction---the coupling matrix is
physically symmetric ($\kappa_{ji} = \kappa_{ij}$), so $W^T$ is
automatically implemented.

\paragraph{Backward pass time.}
\begin{equation}
\tau_{\mathrm{backward}} = \tau_{\mathrm{prop}} + \tau_{\mathrm{settle}}
\approx \tau_{\mathrm{forward}} \approx 1\text{--}2\,\mathrm{ns},
\end{equation}
again independent of layer size.

\subsection{Weight Update: Gradient Descent on $\psi$-Field}\label{sec:weight_update}

The weight update rule is:
\begin{equation}
w_{ji}^{(t+1)} = w_{ji}^{(t)} - \eta\,\frac{\partial\mathcal{L}}{\partial w_{ji}}
= w_{ji}^{(t)} - \eta\,\delta_j\,x_i.
\label{eq:sgd}
\end{equation}

\paragraph{Physical implementation.}
The coupling strength $\kappa_{ji}$ is modified by adjusting the
intermediary cavity detuning $\delta\omega_{ji}$ (fast, electronic) or
the aperture geometry (slow, piezoelectric):
\begin{equation}
\kappa_{ji}^{\mathrm{eff}} = \frac{\kappa_0^2}{\delta\omega_{ji}^2 + \gamma_{\mathrm{med}}^2},
\end{equation}
so
\begin{equation}
\Delta w_{ji} = -\frac{2\kappa_0^2\,\delta\omega_{ji}}{(\delta\omega_{ji}^2 + \gamma_{\mathrm{med}}^2)^2}\,\Delta(\delta\omega_{ji}).
\end{equation}

Setting $\Delta w_{ji} = -\eta\,\delta_j\,x_i$ determines the required
detuning change:
\begin{equation}
\Delta(\delta\omega_{ji}) = \frac{\eta\,\delta_j\,x_i\,(\delta\omega_{ji}^2 + \gamma_{\mathrm{med}}^2)^2}{2\kappa_0^2\,\delta\omega_{ji}}.
\end{equation}

\subsection{TOPS/W Analysis}\label{sec:topsw}

\paragraph{Operations per second.}
For a single layer with $N_{\mathrm{in}} = N_{\mathrm{out}} = N$
neurons, operating at clock frequency $f$:
\begin{equation}
\mathrm{OPS} = N^2 \times f
\end{equation}
(each forward pass performs $N^2$ MAC operations).

\paragraph{Power consumption (computation only).}
\begin{equation}
P_{\mathrm{compute}} = N \times \frac{\pi\alpha_{\mathrm{rel}}\kB T}{Q} \times f.
\end{equation}

\paragraph{TOPS/W (computation).}
\begin{equation}
\mathrm{TOPS/W}_{\mathrm{compute}} = \frac{N^2 f \times 10^{-12}}{P_{\mathrm{compute}}}
= \frac{N \times 10^{-12}\,Q}{\pi\alpha_{\mathrm{rel}}\kB T}.
\label{eq:topsw_compute}
\end{equation}

For $N = 1024$, $Q = 10^{10}$, $T = 1.5\,\mathrm{K}$:
\begin{equation}
\mathrm{TOPS/W}_{\mathrm{compute}} = \frac{1024 \times 10^{-12} \times 10^{10}}{\pi \times 40 \times 1.38\times 10^{-23} \times 1.5}
\approx 3.9\times 10^{27}\,\mathrm{TOPS/W}.
\end{equation}

\paragraph{System-level TOPS/W (including cryogenics).}
Cryogenic cooling at 1.5\,K requires approximately 1000\,W of wall-plug
power per 1\,W of cooling at the cold stage (Carnot efficiency $\times$
practical efficiency: $\eta_{\mathrm{Carnot}} = T_c/(T_h - T_c) \approx
0.005$, practical COP $\approx 10^{-3}$).

The cryogenic overhead dominates. For a system cooling budget of
$P_{\mathrm{cryo,cold}} = 1\,\mathrm{W}$ (supporting many thousands of
cavities), the wall power is $P_{\mathrm{wall}} \approx 1\,\mathrm{kW}$.

The compute power for the neural layer is $P_{\mathrm{compute}} \approx
10^{-19}\,\mathrm{W}$---completely negligible compared to cryogenics. We
can therefore pack as many layers as the cryogenic budget allows.

For a system with $M$ layers of size $N = 1024$:
\begin{equation}
\mathrm{TOPS/W}_{\mathrm{system}} = \frac{M \times N^2 \times f \times 10^{-12}}{P_{\mathrm{wall}}}
= \frac{M \times 1024^2 \times 10^9 \times 10^{-12}}{10^3}.
\end{equation}

For $M = 1000$ layers (feasible within cryogenic budget):
\begin{equation}
\mathrm{TOPS/W}_{\mathrm{system}} \approx \frac{1000 \times 10^6 \times 10^{-3}}{10^3}
= 10^3\,\mathrm{TOPS/W}.
\end{equation}

For $M = 10^6$ layers (scaling with miniaturized cavities at 100\,GHz):
\begin{equation}
\mathrm{TOPS/W}_{\mathrm{system}} \approx 10^6\,\mathrm{TOPS/W}.
\end{equation}

\subsection{Comparison with State-of-the-Art AI Accelerators}
\label{sec:ai_comparison}

\begin{table}[h]
\centering
\caption{AI accelerator comparison: theoretical peak TOPS/W.}
\label{tab:ai_comparison}
\begin{tabular}{@{}lccc@{}}
\toprule
\textbf{Platform} & \textbf{TOPS/W} & \textbf{Precision} & \textbf{$T$ (K)} \\
\midrule
NVIDIA H100 (SXM5) & $\sim 2\times 10^3$ & FP8 & 300 \\
Google TPU v5e & $\sim 3\times 10^3$ & BF16/INT8 & 300 \\
Intel Loihi 2 & $\sim 10^4$ & INT8 spike & 300 \\
Cerebras CS-3 & $\sim 10^3$ & FP16 & 300 \\
Groq LPU & $\sim 3\times 10^3$ & INT8 & 300 \\
\midrule
$\psi$-field (1k layers) & $\sim 10^3$ & Analog & 1.5 \\
$\psi$-field ($10^6$ layers) & $\sim 10^6$ & Analog & 1.5 \\
$\psi$-field (compute only) & $\sim 10^{27}$ & Analog & 1.5 \\
\bottomrule
\end{tabular}
\end{table}

\paragraph{Key observations.}
\begin{enumerate}
\item At modest scale ($\sim 1000$ cavity layers), $\psi$-field neural
processing achieves system-level TOPS/W comparable to current GPU
accelerators, despite the cryogenic overhead.
\item At large scale ($> 10^5$ layers, requiring miniaturized cavities),
$\psi$-field systems significantly exceed all room-temperature
accelerators.
\item The compute-only TOPS/W of $\sim 10^{27}$ represents the
\emph{intrinsic} efficiency of $\psi$-field physics---the
computation dissipates effectively zero energy. All practical limits come
from the cryogenic infrastructure.
\item The analog precision of $\psi$-field computation
(continuous-valued weights and activations) avoids quantization errors
that limit digital accelerators to FP8/INT8 formats.
\end{enumerate}

\subsection{Mapping Specific Network Architectures}\label{sec:architectures}

\subsubsection{Convolutional Layer}

A convolutional layer with kernel size $k \times k$, $C_{\mathrm{in}}$
input channels, and $C_{\mathrm{out}}$ output channels requires:
\begin{equation}
N_{\mathrm{MAC}} = k^2 \times C_{\mathrm{in}} \times C_{\mathrm{out}} \times H \times W
\end{equation}
MACs per forward pass, where $H \times W$ is the spatial resolution.

In $\psi$-field implementation, weight sharing is achieved by using the
\emph{same} coupling matrix for all spatial positions, with the input
shifted through the coupling apertures sequentially. The computation
time per spatial position is $\tau_{\mathrm{forward}}$, so the total
convolution time is:
\begin{equation}
\tau_{\mathrm{conv}} = H \times W \times \tau_{\mathrm{forward}} \approx H \times W \times 2\,\mathrm{ns}.
\end{equation}

\subsubsection{Attention Layer (Transformer)}

Self-attention requires computing $QK^T$ and softmax. In $\psi$-field:
\begin{itemize}
\item \textbf{$QK^T$}: Matrix multiplication performed by coupling between
query and key cavity arrays. Time: $\tau_{\mathrm{forward}}$ (independent
of sequence length).
\item \textbf{Softmax}: Implemented by mutual inhibition between output
cavities. The competition dynamics of coupled bistable systems naturally
produce a softmax-like winner-take-all output.
\item \textbf{Value weighting}: Second coupling matrix multiplies
attention weights by value vectors.
\end{itemize}

Total attention layer time: $\sim 3\tau_{\mathrm{forward}} \approx
5\,\mathrm{ns}$.

%=============================================================================
\section{Discussion}\label{sec:discussion}
%=============================================================================

\subsection{Thermodynamic Consistency}

The sub-Landauer result (Theorem~\ref{thm:entropy_production}) is
thermodynamically consistent because:
\begin{enumerate}
\item It applies only to \emph{reversible} logic operations that erase
no information.
\item The dissipation $\pi\alpha_{\mathrm{rel}}\kB T/Q > 0$ is strictly
positive, satisfying the second law.
\item The Jarzynski equality~\eqref{eq:jarzynski} and Crooks
theorem~\eqref{eq:crooks} are exactly satisfied.
\item For irreversible operations, the full Landauer cost
$\kB T\ln 2$ is paid in addition to the friction loss
(Eq.~\ref{eq:irrev_dissipation}).
\end{enumerate}

The crucial physical point is that high-$Q$ superconducting cavities
provide an environment where friction is extraordinarily weak
($\gamma/\omega_0 \sim 10^{-10}$), allowing the system to evolve
nearly-Hamiltonianly. This is not achievable in any electronic system
due to the fundamental resistivity of conductors (even superconductors
have AC losses that limit $Q$ in Josephson circuits to $\sim 10^6$).

\subsection{Scalability Challenges}

The primary bottleneck for $\psi$-field computing is not energy
efficiency but \emph{physical size}. At 1.3\,GHz, each cavity is
$\sim 20\,\mathrm{cm}$. A 1024-neuron layer requires $\sim 1000$
cavities, occupying $\sim 10\,\mathrm{m}^3$ at this frequency.

Scaling to higher frequencies reduces cavity size: at $f = 100\,\mathrm{GHz}$,
cavities are $\sim 1.5\,\mathrm{mm}$, and a 1024-neuron layer fits in
$\sim 10\,\mathrm{cm}^3$---comparable to a modern chip package.

The challenge is maintaining high $Q$ at higher frequencies. Current
SRF technology achieves $Q \sim 10^{10}$ at 1.3\,GHz, $Q \sim 10^9$ at
10\,GHz, and $Q \sim 10^7$--$10^8$ at 100\,GHz. Even at $Q = 10^7$, the
sub-Landauer condition ($Q > 182$) is satisfied by a factor of $5\times 10^4$.

\subsection{Experimental Verification Path}

The predictions of this paper can be tested experimentally:
\begin{enumerate}
\item \textbf{Sub-Landauer dissipation}: Measure heat production per
switching event in an SRF cavity with Josephson nonlinearity, using
calorimetric or noise thermometry techniques. The predicted dissipation
of $\sim 10^{-31}\,\mathrm{J}$ at $Q = 10^{10}$ is below current
measurement sensitivity ($\sim 10^{-23}\,\mathrm{J}$), but the scaling
$E_{\mathrm{diss}} \propto 1/Q$ can be verified by operating at
deliberately reduced $Q$ values ($Q = 10^3$--$10^6$).
\item \textbf{Radiation hardness}: Irradiate a $\psi$-field memory cell
with known particle beams and measure the SEU rate as a function of the
number of participating modes.
\item \textbf{Neural layer}: Implement a small ($4 \times 4$) coupled
cavity array and demonstrate matrix-vector multiplication followed by
learning.
\end{enumerate}

%=============================================================================
\section{Conclusions}\label{sec:conclusions}
%=============================================================================

We have performed a deep quantitative analysis of $\psi$-field computing
across five dimensions:

\begin{enumerate}
\item \textbf{Sub-Landauer thermodynamics}: Starting from the Jarzynski
equality, we proved rigorously that reversible $\psi$-field logic with
$Q > 182$ beats the Landauer bound, with exact entropy production
$\Delta S_{\mathrm{op}} = \pi\alpha_{\mathrm{rel}}\kB/Q$ per operation.
At $Q = 10^{10}$, this is $5.5\times 10^7$ times below the Landauer
limit---without violating the second law.

\item \textbf{Literature comparison}: Quantitative benchmarking against
12 technologies (B\'er\'ut experiment, RSFQ, AQFP, adiabatic CMOS,
photonic, neuromorphic, memristive) shows $\psi$-field computing holds a
$10^{10}$--$10^{19}\times$ energy advantage per operation over all
existing platforms.

\item \textbf{4-bit ALU}: A complete reversible ALU requires 47 cavities
and performs addition at $\sim 10^{-29}\,\mathrm{J}$ per operation---$10^{13}\times$
more efficient than 7\,nm CMOS and $10^{9}\times$ more efficient than
AQFP.

\item \textbf{Radiation hardness}: Multi-mode encoding provides
exponential SEU suppression. With 5 participating modes,
$\sigma_{\mathrm{SEU}} < 10^{-19}\,\mathrm{cm}^2/\mathrm{bit}$,
surpassing all existing memory technologies.

\item \textbf{Neural processing}: A $\psi$-field neural layer performs
forward and backward passes in $O(1)$ time independent of layer size.
System-level TOPS/W of $10^3$--$10^6$ is achievable, matching or
exceeding NVIDIA H100 and Google TPU v5 while providing analog precision.
\end{enumerate}

The $\psi$-field computing paradigm represents a fundamentally new point
in the computing landscape: computation at the thermodynamic limit of
what physics allows, enabled by the extraordinary coherence properties
of superconducting cavities operating on the DFD scalar refractive field.

%=============================================================================
% BIBLIOGRAPHY
%=============================================================================
\begin{thebibliography}{99}

\bibitem{Alcock2025psiComputing}
G.~T.~Alcock, ``$\psi$-Field Computing: Logic, Memory, and Neural
Processing in Scalar Refractive Geometry,'' DFD Research Programme (2025).

\bibitem{Alcock2025DFD}
G.~T.~Alcock, \textit{Density Field Dynamics: A Complete Unified Theory},
v3.2 (2025).

\bibitem{Landauer1961}
R.~Landauer, ``Irreversibility and Heat Generation in the Computing
Process,'' \textit{IBM J. Res. Dev.} \textbf{5}, 183 (1961).

\bibitem{Bennett1973}
C.~H.~Bennett, ``Logical Reversibility of Computation,''
\textit{IBM J. Res. Dev.} \textbf{17}, 525 (1973).

\bibitem{Bennett1982}
C.~H.~Bennett, ``The Thermodynamics of Computation---a Review,''
\textit{Int. J. Theor. Phys.} \textbf{21}, 905 (1982).

\bibitem{Jarzynski1997}
C.~Jarzynski, ``Nonequilibrium Equality for Free Energy Differences,''
\textit{Phys. Rev. Lett.} \textbf{78}, 2690 (1997).

\bibitem{Crooks1999}
G.~E.~Crooks, ``Entropy Production Fluctuation Theorem and the
Nonequilibrium Work Relation for Free Energy Differences,''
\textit{Phys. Rev. E} \textbf{60}, 2721 (1999).

\bibitem{Berut2012}
A.~B\'er\'ut, A.~Arakelyan, A.~Petrosyan, S.~Ciliberto, R.~Dillenschneider,
and E.~Lutz, ``Experimental Verification of Landauer's Principle Linking
Information and Thermodynamics,'' \textit{Nature} \textbf{483}, 187 (2012).

\bibitem{Jun2014}
Y.~Jun, M.~Gavrilov, and J.~Bhcek, ``High-Precision Test of Landauer's
Principle in a Feedback Trap,'' \textit{Phys. Rev. Lett.} \textbf{113},
190601 (2014).

\bibitem{Hong2016}
J.~Hong, B.~Lambson, S.~Dhuey, and J.~Bokor, ``Experimental Test of
Landauer's Principle in Single-Bit Operations on Nanomagnetic Memory
Bits,'' \textit{Sci. Adv.} \textbf{2}, e1501492 (2016).

\bibitem{Gaudenzi2018}
R.~Gaudenzi, E.~Burzur\'{i}, S.~Maegawa, H.~S.~J.~van~der~Zant, and
F.~Luis, ``Quantum Landauer Erasure with a Molecular Nanomagnet,''
\textit{Nat. Phys.} \textbf{14}, 565 (2018).

\bibitem{FredkinToffoli1982}
E.~Fredkin and T.~Toffoli, ``Conservative Logic,''
\textit{Int. J. Theor. Phys.} \textbf{21}, 219 (1982).

\bibitem{Athas1994}
W.~C.~Athas, L.~J.~Svensson, J.~G.~Koller, N.~Tzartzanis, and
E.~Y.-C.~Chou, ``Low-Power Digital Systems Based on Adiabatic-Switching
Principles,'' \textit{IEEE Trans. VLSI Syst.} \textbf{2}, 398 (1994).

\bibitem{Younis1994}
S.~G.~Younis and T.~F.~Knight, ``Asymptotically Zero Energy Split-Level
Charge Recovery Logic,'' in \textit{Proc. ISLPED} (1994), pp.~177--182.

\bibitem{Likharev1991}
K.~K.~Likharev and V.~K.~Semenov, ``RSFQ Logic/Memory Family: A New
Josephson-Junction Technology for Sub-Terahertz-Clock-Frequency Digital
Systems,'' \textit{IEEE Trans. Appl. Supercond.} \textbf{1}, 3 (1991).

\bibitem{Takeuchi2013}
N.~Takeuchi, D.~Ozawa, Y.~Yamanashi, and N.~Yoshikawa, ``An Adiabatic
Quantum Flux Parametron as an Ultra-Low-Power Logic Device,''
\textit{Supercond. Sci. Technol.} \textbf{26}, 035010 (2013).

\bibitem{Takeuchi2017}
N.~Takeuchi, T.~Yamae, C.~L.~Ayala, H.~Suzuki, and N.~Yoshikawa,
``Adiabatic Quantum-Flux-Parametron: A Tutorial Review,''
\textit{IEICE Trans. Electron.} \textbf{E101.C}, 364 (2018).

\bibitem{Shen2017}
Y.~Shen, N.~C.~Harris, S.~Skirlo, \textit{et al.}, ``Deep Learning
with Coherent Nanophotonic Circuits,'' \textit{Nat. Photonics}
\textbf{11}, 441 (2017).

\bibitem{Harris2018}
N.~C.~Harris, G.~R.~Steinbrecher, M.~Prabhu, \textit{et al.},
``Linear Programmable Nanophotonic Processors,'' \textit{Optica}
\textbf{5}, 1623 (2018).

\bibitem{Larger2012}
L.~Larger, M.~C.~Soriano, D.~Brunner, \textit{et al.}, ``Photonic
Information Processing Beyond Turing: an Optoelectronic Implementation
of Reservoir Computing,'' \textit{Opt. Express} \textbf{20}, 3241 (2012).

\bibitem{Davies2018}
M.~Davies, N.~Srinivasa, T.-H.~Lin, \textit{et al.}, ``Loihi: A
Neuromorphic Manycore Processor with On-Chip Learning,''
\textit{IEEE Micro} \textbf{38}, 82 (2018).

\bibitem{Orchard2021}
G.~Orchard, E.~P.~Frady, D.~B.~D.~Rubin, S.~Sanborn, S.~B.~Shrestha,
F.~T.~Sommer, and M.~Davies, ``Efficient Neuromorphic Signal Processing
with Loihi 2,'' in \textit{Proc. IEEE Workshop on Signal Processing
Systems} (2021).

\bibitem{Merolla2014}
P.~A.~Merolla, J.~V.~Arthur, R.~Alvarez-Icaza, \textit{et al.},
``A Million Spiking-Neuron Integrated Circuit with a Scalable
Communication Network and Interface,'' \textit{Science} \textbf{345},
668 (2014).

\bibitem{Prezioso2015}
M.~Prezioso, F.~Merrikh-Bayat, B.~D.~Hoskins, G.~C.~Adam,
K.~K.~Likharev, and D.~B.~Strukov, ``Training and Operation of an
Integrated Neuromorphic Network Based on Metal-Oxide Memristors,''
\textit{Nature} \textbf{521}, 61 (2015).

\bibitem{Li2018Nature}
C.~Li, D.~Belkin, Y.~Li, \textit{et al.}, ``Efficient and
Self-Adaptive In-Situ Learning in Multilayer Memristor Neural Networks,''
\textit{Nat. Commun.} \textbf{9}, 2385 (2018).

\bibitem{Cuccaro2004}
S.~A.~Cuccaro, T.~G.~Draper, S.~A.~Kutin, and D.~P.~Moulton,
``A New Quantum Ripple-Carry Addition Circuit,'' arXiv:quant-ph/0410184
(2004).

\end{thebibliography}

\end{document}
